\chapter{El dinero}
\label{ch:money}

\section{Qué hizo mal el Ordenamiento}

El capítulo~\ref{ch:crisis} describió la unificación monetaria de enero de 2021. Este
capítulo está escrito contra ella, así que conviene reformular el fallo con
precisión, porque la lección es sobre la secuencia y no sobre la dirección.

La unificación era correcta y llevaba veinte años siéndolo. Con la tasa de los hogares
separada veinticinco veces de la tasa empresarial, ninguna contabilidad de empresa
cubana significaba nada, y una economía cuyas cuentas son ficción no puede reformarla
ni siquiera quien quiera reformarla. Lo que salió mal fue todo lo que la rodeaba.

Se hizo sin reservas, y una tasa fija es la promesa de vender dólares a esa tasa. Cuba
no tenía dólares, de modo que la promesa no se cumplió y en semanas reabrió el mercado
paralelo ---y una vez que existe una tasa paralela, la tasa oficial deja de ser un
precio y se convierte en un derecho de asignación, repartido por quien controle el
acceso a ella---.

Se hizo sin respuesta de la oferta. Los salarios y las pensiones se multiplicaron
varias veces en una economía donde la cantidad de bienes estaba cayendo. Eso no es una
decisión distributiva; es la definición de una inflación.

Se hizo con el sector privado todavía cerrado. La congelación de licencias de 2017
seguía vigente y la ley que dio personalidad jurídica a las empresas privadas estaba a
ocho meses, de modo que a la única parte de la economía capaz de responder a precios
más altos con más producción se le impedía legalmente hacerlo.

Y se hizo en el fondo del peor choque externo desde 1993, cuando el orden correcto es
estabilizar primero y unificar después.

La comparación más dura es interna. El mismo Estado, en 1994--96, con menos recursos y
una posición de partida peor, recortó el déficit de alrededor del treinta por ciento
del producto al dos, llevó la tasa informal de ciento cincuenta a veinticinco, y
estabilizó el peso sin asistencia externa ---y lo hizo \emph{antes} que cualquier otra
cosa---. El capítulo~\ref{ch:dollarreform} cuenta esa historia. Cuba ha demostrado que
sabe hacer esto en el orden correcto. En 2021 no lo hizo.

\section{El régimen monetario}

El libro toma aquí una posición, y es la recomendación más discutible de la Parte III.

\emph{Cuba debe dolarizarse oficial y unilateralmente, como régimen de entrada, con una
salida declarada.}

El argumento no es que la dolarización sea óptima en general. No lo es. Es que las
alternativas exigen todas algo que Cuba no tiene.

Una \emph{flotación administrada con banco central independiente} exige un banco
central cuyos compromisos se crean. El de Cuba ha destruido los ahorros de sus
ciudadanos dos veces en treinta años ---en 1993--94 y otra vez en 2021--- y su
capacidad de prometer algo sobre el valor del dinero es, como cuestión de hecho
observado, nula. Una credibilidad así se reconstruye a lo largo de una década de
conducta demostrada, y el diseño necesita una moneda que funcione en el mes uno.

Una \emph{caja de conversión}, al modo báltico o búlgaro, conserva una moneda nacional
plenamente respaldada por reservas. Es la respuesta correcta para un país que tiene
reservas. Cuba no tiene, no puede pedirlas prestadas, y tendría que adquirirlas antes
de que la caja pudiera abrir ---lo que es un rodeo de cinco años en el mejor caso---.

Un \emph{deslizamiento controlado} exige reservas que defender y discrecionalidad para
ajustar, y la discrecionalidad sobre un tipo de cambio en un país con la historia
institucional de Cuba es una renta, no una política. El capítulo~\ref{ch:integrity}
señalará que la mayor superficie individual de corrupción de la economía cubana durante
treinta años ha sido el acceso a divisa a la tasa oficial.

La dolarización es el régimen que no exige reservas, ni credibilidad, ni
discrecionalidad, que es exactamente la combinación que especifica el conjunto de
restricciones cubano. Es además hacia donde el país ya va: los hogares fijan precios en
dólares, las remesas llegan en dólares, y desde 2024 las propias tiendas del Estado
venden en dólares. La propuesta es volver oficial y universal lo que ya es parcial e
inequitativo ---porque una economía parcialmente dolarizada es el peor de todos los
arreglos, ya que divide a la población entre quienes tienen acceso a divisa y quienes
cobran en otra cosa, que es precisamente la inversión que describía el
capítulo~\ref{ch:dollarreform} y que está ocurriendo otra vez---.

Los costos son reales y hay que declararlos. Cuba renuncia a la política monetaria,
renuncia al tipo de cambio como mecanismo de ajuste, renuncia al señoreaje, y renuncia
a un prestamista de última instancia. Los dos primeros importan menos que en otros
lugares: los exportadores cubanos ---turismo, servicios profesionales, níquel, y los
servicios remotos del capítulo~\ref{ch:talent}--- ya fijan precios en dólares, de modo
que una devaluación no compraría la competitividad que le compra a un exportador
manufacturero. El señoreaje sobre una moneda que nadie quiere tener es pequeño. El
prestamista de última instancia es la pérdida genuina, y muerde con más fuerza
exactamente en el escenario que Cuba enfrenta con regularidad ---un huracán de
categoría 5---, razón por la cual el fondo de estabilización del
Diseño~\ref{d:rentrule} y la línea de crédito contingente de más abajo no son extras
opcionales sino la otra mitad de esta elección de régimen.

Ecuador, El Salvador y Panamá son los comparadores. El caso ecuatoriano es el
instructivo: la dolarización se adoptó en 2000 en medio de un derrumbe bancario, la
mayoría de los economistas esperaba que fracasara, y ha sobrevivido veinticinco años y
varios gobiernos a los que no les gustaba ---porque resultó ser enormemente popular
entre gente que había perdido sus ahorros, y ningún gobierno ha podido revertirla---.
Eso es un dispositivo de compromiso en el sentido del capítulo~\ref{ch:polity}: una
política que se vuelve difícil de deshacer porque la población no lo permitiría.

\begin{design}{El régimen monetario}
\label{d:regime}
\begin{rules}
\rl{1}{Adopción} El dólar estadounidense se adopta como moneda de curso legal y unidad
de cuenta. El peso cubano se retira según un calendario publicado a una tasa única de
conversión aplicada simultáneamente a todos los saldos, salarios, precios, contratos y
cuentas públicas.

\rl{2}{La tasa de conversión} Se fija en la tasa de mercado informal vigente o cerca de
ella, no en una tasa administrativa. Fijarla demasiado fuerte es lo que destruye los
ahorros; fijarla en la tasa de mercado simplemente reconoce una pérdida que ya ocurrió,
que es la diferencia entre una reforma y una confiscación.

\rl{3}{Tasa única, de inmediato} No hay tasa dual transitoria, ni tasa preferencial para
las empresas estatales, ni asignación de divisa a tasa privilegiada a nadie, para nada,
incluidas las importaciones de alimentos y combustible. Todos los arreglos de ese tipo
en la historia cubana han terminado siendo una renta.

\rl{4}{Protección de salarios y pensiones} Los salarios y pensiones públicos se
convierten a la misma tasa y después se recalibran según el Diseño~\ref{d:pay}, con un
piso que garantice que ningún pensionado vea caer su ingreso real en la conversión. Es
la cláusula que vuelve sobrevivible el régimen, y está costeada en el
capítulo~\ref{ch:sequence}.

\rl{5}{La salida} Un conjunto publicado de condiciones ---reservas, historial fiscal,
inflación, una autoridad supervisora establecida--- bajo las cuales pueda reintroducirse
más adelante una moneda nacional, junto con el procedimiento. Nombrar la salida es lo
que distingue una elección meditada de una rendición, y escribirla no cuesta nada.
\end{rules}
\end{design}

\begin{design}{Qué sustituye al banco central}
\label{d:bank}
\begin{rules}
\rl{1}{Función} Bajo dolarización el banco central no emite dinero. Pasa a ser
supervisor bancario, autoridad de pagos y gestor del fondo de estabilización, con el
mandato, el presupuesto y las obligaciones de publicación del Diseño~\ref{d:courts}.

\rl{2}{Sin financiamiento fiscal} La prohibición de prestar al gobierno es absoluta y
constitucional. Es la cláusula monetaria más importante del libro, porque el
financiamiento monetario del presupuesto es como Cuba pagó su Estado de bienestar en
los años en que ni Moscú ni Caracas estaban pagando, y es lo que el
capítulo~\ref{ch:welfare} está diseñado para sustituir.

\rl{3}{Liquidez} Una facilidad de liquidez financiada con el fondo de estabilización y,
si puede obtenerse, una línea de crédito contingente con un socio multilateral o
bilateral ---el sustituto del prestamista de última instancia al que renuncia la
cláusula~1---. Obtenerla es un objetivo de primer orden de la normalización de la deuda
que viene abajo, no una idea de última hora.

\rl{4}{Supervisión bancaria} Estándares de capital, liquidez y divulgación conforme a
normas internacionales, aplicados por igual a bancos estatales y privados, con un
régimen de resolución que exista antes de necesitarse.
\end{rules}
\end{design}

\section{Deuda, atrasos y reclamaciones}

Cuba le debe dinero a casi todo el mundo y no le ha pagado a la mayoría. El acervo son
varios problemas distintos que convencionalmente se tratan por separado y que este
libro propone tratar en una sola operación.

\emph{Deuda oficial del Club de París.} El arreglo de 2015 condonó buena parte de los
atrasos acumulados y reprogramó el resto a dieciocho años. Cuba empezó a fallar pagos en
tres años y el acuerdo está en incumplimiento.

\emph{Reclamaciones rusas.} El noventa por ciento de la deuda de la era soviética se
condonó en 2014; el resto se convirtió en compromisos de inversión, en gran medida no
materializados.

\emph{Reclamaciones chinas y otras bilaterales.} Sustancialmente no divulgadas, lo que
es en sí mismo un problema: una reestructuración no puede negociarse contra un pasivo
desconocido, y el primer acto de cualquier proceso es publicar el acervo.

\emph{Atrasos con proveedores.} Facturas impagas a socios comerciales extranjeros, en
algunos casos hasta 2009. Comercialmente son los más dañinos, porque los acreedores son
las mismas empresas que de otro modo suministrarían las importaciones que necesita la
recuperación.

\emph{Cuentas congeladas.} Saldos de empresas extranjeras bloqueados desde 2009.
Jurídicamente no son deuda; comercialmente son la razón de que nadie confíe en una
cuenta cubana.

\emph{Reclamaciones por expropiación.} La Comisión de Resolución de Reclamaciones
Extranjeras registró 8.821 reportes de bienes incautados y certificó 5.913 de ellos, por
un principal de 1.900 millones de dólares. Al 6 por ciento de interés simple que
especifica el estatuto que la rige, eso está hoy valorado en algo más de 9.000 millones
\citep{fcsc}. Se estima que solo unas 913 de las reclamaciones certificadas ---cerca del
15 por ciento--- cumplen las condiciones para una demanda bajo el Título III. Junto a
ellas hay un cuerpo mucho mayor y no certificado de reclamaciones de cubanoamericanos y
de cubanos que nunca se fueron.

\begin{design}{Una reestructuración, no seis}
\label{d:debt}
\begin{rules}
\rl{1}{Publicar primero} Una declaración pública completa y auditada de cada obligación
externa ---oficial, bilateral, comercial, atrasos, saldos congelados y contingentes---
antes de que se abra cualquier negociación. Ninguna reestructuración ha tenido éxito
jamás contra un acervo no divulgado, y la divulgación es además el primer acto del
régimen de transparencia del capítulo~\ref{ch:integrity}.

\rl{2}{Un solo proceso} Una única negociación de trato comparable que abarque todas las
clases de reclamación, las de expropiación incluidas. Saldar la deuda soberana dejando
pendientes las reclamaciones por expropiación deja a todo inversionista extranjero
expuesto a litigio bajo el Título III y por tanto deja el país no invertible; saldar las
reclamaciones fuera de una reestructuración general le da prioridad a una clase de
acreedor y hunde al resto.

\rl{3}{Instrumentos} Bonos de largo plazo y cupón bajo con un componente ligado al
crecimiento, de modo que los acreedores participen en una recuperación que están
ayudando a financiar y el calendario no apriete en los años en que el plan necesita
caja. No una quita disfrazada de extensión de plazos: declarar la reducción y
negociarla.

\rl{4}{Reclamaciones saldadas en papel, dentro de un tope} Según el
capítulo~\ref{ch:capital}, las reclamaciones por expropiación se saldan con
compensación y no con restitución, en los mismos instrumentos, dentro de un tope
agregado anunciado de antemano. Ningún cubano es sacado de una vivienda.

\rl{5}{Membresía multilateral} El ingreso al Fondo Monetario Internacional y a los bancos
multilaterales de desarrollo se persigue en paralelo y es un objetivo de la
reestructuración antes que una consecuencia de ella, porque su asistencia técnica y su
financiamiento contingente son lo que exige la cláusula~3 del Diseño~\ref{d:bank}.

\rl{6}{La conducta antes que los términos} Los saldos congelados del
capítulo~\ref{ch:venezuela} se liberan y las facturas corrientes de proveedores se pagan
puntualmente desde el primer mes, antes de y con independencia de cualquier negociación.
Un historial crediticio se repara con conducta a lo largo de años y no con un acuerdo
firmado, y esta es la credibilidad más barata disponible.
\end{rules}
\end{design}

\section{El presupuesto}

El capítulo~\ref{ch:welfare} es dueño de qué se grava. Este capítulo es dueño del
déficit, de cómo se financia, y del ancla que impide que el arreglo se deshaga.

\begin{design}{Arquitectura fiscal}
\label{d:fiscal}
\begin{rules}
\rl{1}{Un solo presupuesto} Todos los ingresos y gastos del Estado, incluidos los de las
empresas estatales y los de los conglomerados vinculados a los militares, aparecen en un
único presupuesto publicado según estándares internacionales de clasificación. Hoy una
parte grande de la hacienda pública cubana está fuera de presupuesto, lo que significa
que nadie ---el gobierno incluido--- conoce la posición fiscal.

\rl{2}{Sin financiamiento monetario} Según el Diseño~\ref{d:bank}.2, y bajo dolarización
es además mecánicamente imposible y no solo prohibido, que es uno de los principales
atractivos del régimen.

\rl{3}{Ancla} La regla de equilibrio corriente del Diseño~\ref{d:rentrule}.3, con un
consejo fiscal independiente que informa anualmente a la asamblea y publica sus propias
proyecciones.

\rl{4}{Presupuestos duros para las empresas} Las empresas estatales enfrentan una
restricción presupuestaria dura con un procedimiento concursal publicado. El subsidio a
una empresa en pérdidas, allí donde continúe, es una partida presupuestaria explícita y
visible, y no un sobregiro bancario que nunca se reclama.

\rl{5}{Secuenciado} El presupuesto se lleva al equilibrio antes de que el sistema
tributario sea capaz de financiar el compromiso de bienestar, lo que significa que la
primera fase corre sobre disciplina de gasto y la regla de la renta, y que la
convergencia hacia los niveles de prestación del capítulo~\ref{ch:welfare} espera por la
base. Decir cuál de esas cosas va primero es todo el contenido de la palabra
``secuencia''.
\end{rules}
\end{design}

\section{La libreta}

La libreta se introdujo en marzo de 1962 como medida temporal y tiene sesenta y cuatro
años. Es un subsidio universal en especie, mal focalizado ---el ministro y el
pensionado reciben lo mismo---, caro en divisas y cada vez menos entregado. Es también,
para un número sustancial de hogares de la tercera edad y del oriente, la diferencia
entre comer y no comer.

La transición es un problema genuino de diseño y no una cuestión de voluntad política, y
hacerlo mal mata de hambre a la gente.

\begin{design}{De la ración al ingreso}
\label{d:libreta}
\begin{rules}
\rl{1}{Nada se retira antes de que funcione el sustituto} Las prestaciones universales en
efectivo del Diseño~\ref{d:benefit} tienen que estar pagándose, íntegras, con
fiabilidad, a través de un sistema de pagos que llegue a cada municipio, durante un
mínimo de dos años consecutivos antes de que se retire ningún renglón de la libreta.

\rl{2}{Sustituir, no focalizar} La libreta se convierte en efectivo pagado a todos los
hogares, no en una prestación focalizada. La universalidad se conserva por la razón del
capítulo~\ref{ch:welfare}: una comprobación de medios es una decisión discrecional y
Cuba no puede administrar una honestamente.

\rl{3}{Indexado} El monto en efectivo se indexa al precio publicado de una canasta, de
modo que la conversión no sea una transferencia única cuyo valor elimine la inflación
---que es lo que le ocurrió a los aumentos salariales de enero de 2021---.

\rl{4}{Secuencia por renglón} Los renglones se convierten uno a uno, empezando por
aquellos cuyo abasto de mercado sea suficiente y terminando por aquellos cuyo abasto
dependa del éxito del capítulo~\ref{ch:land}. El arroz, el aceite y la leche van al
final.

\rl{5}{Reversible} Si el sustituto falla ---si el abasto se derrumba o el sistema de
pagos no llega a una provincia--- la ración se restablece para ese renglón y esa
provincia, automáticamente, por un disparador publicado. Es el único lugar de la Parte
III donde la reversibilidad es una virtud.
\end{rules}
\end{design}

\section{Los pagos}

La sección menos vistosa de la Parte III y de las más vinculantes.

Un turista no puede pagar en un restaurante cubano con una tarjeta corriente. Un
trabajador remoto en La Habana no puede recibir dinero de un cliente en Madrid en una
cuenta que controle. Una empresa privada no puede pagarle a un proveedor extranjero. Un
pensionado en Baracoa no puede recibir una transferencia sin una oficina física. La
estrategia de valor por visitante del capítulo~\ref{ch:tourism}, la proposición entera
del capítulo~\ref{ch:talent} y la conversión del Diseño~\ref{d:libreta} fracasan todas
sin rieles de pago, y ninguna de ellas puede construirse hasta que existan los rieles.

\begin{design}{Rieles de pago}
\label{d:pay2}
\begin{rules}
\rl{1}{Banca corresponsal} Restablecer las relaciones de corresponsalía con bancos
internacionales es un objetivo de primer orden y la principal razón práctica de la
cláusula~6 del Diseño~\ref{d:debt} y de los estándares antilavado de la cláusula~4 de
abajo. Bajo el embargo es difícil y, para las instituciones no norteamericanas, no
imposible.

\rl{2}{Tarjetas y aceptación} Aceptación de tarjetas en cualquier negocio con licencia,
con la comisión al comercio limitada y el plazo de liquidación fijado en el reglamento.
Un restaurante privado que puede cobrar con tarjeta vale varias veces uno que no puede.

\rl{3}{El derecho a tener y a cobrar} Cualquier residente puede tener una cuenta en
divisa, recibir en ella transferencias internacionales y gastar desde ella sin permiso
administrativo. Es la cláusula de la que depende enteramente el
capítulo~\ref{ch:talent}, y es hoy la restricción vinculante sobre las exportaciones
cubanas de servicios ---más que la velocidad de internet y mucho más que las visas---.

\rl{4}{Estándares} Estándares antilavado y de titularidad real conforme a normas
internacionales, supervisados por la autoridad del Diseño~\ref{d:bank}. No son una
concesión a los extranjeros; son el precio de las relaciones de corresponsalía de la
cláusula~1, y son el mismo registro que el capítulo~\ref{ch:integrity} necesita por
razones domésticas.

\rl{5}{Alcance} Una cuenta transaccional básica disponible para todo residente, incluso
por móvil, en cada municipio. La inclusión financiera es precondición del
Diseño~\ref{d:libreta} y del impuesto sobre la propiedad del Diseño~\ref{d:tax}.4, y sin
ella los dos son administrativamente imposibles en el oriente.
\end{rules}
\end{design}

\section{Los primeros quinientos días}

El orden, comprimido. El capítulo~\ref{ch:sequence} lo sitúa dentro del programa
completo; aquí va solo la secuencia monetaria.

\emph{Días 1--100.} Publicar el acervo completo de deuda y atrasos. Pagar las facturas
corrientes de proveedores. Liberar los saldos congelados. Anunciar la fecha de
dolarización y la tasa de conversión. Legislar la prohibición de financiamiento
monetario y la regla de equilibrio corriente. Abrir la negociación de deuda y la
conversación sobre el ingreso al FMI.

\emph{Días 100--250.} Convertir. Tasa única, todos los saldos a la vez, salarios y
pensiones protegidos en el piso. Empezar la aceptación de tarjetas y el despliegue de la
cuenta básica.

\emph{Días 250--500.} Levantar la administración tributaria ---registro, números de
contribuyente, retención sobre la nómina estatal, el régimen presuntivo para la
microempresa--- sabiendo que rendirá poco durante años. Empezar la construcción del IVA.
Establecer el consejo fiscal y publicar su primer informe. Empezar la conversión de la
libreta por los dos o tres renglones cuyo abasto de mercado sea suficiente.

Lo que deliberadamente \emph{no} está en los primeros quinientos días: subir las
prestaciones a niveles europeos, la privatización masiva, y la retirada de cualquier
renglón de la libreta cuyo sustituto no se esté pagando ya. Cada una de esas cosas es de
una fase posterior, e intentarlas aquí es como fracasa la secuencia.
